%============================================================
\chapter{Bài Học: Lòng Dũng Cảm (Courage)}
\begin{center}
  \textit{\color{truyenblue}Dũng cảm không phải không có sợ hãi; mà là đánh giá điều khác quan trọng hơn sợ hãi.}\\
  \textit{(Courage is not the absence of fear; it is deciding something else is more important than fear.)}\\[4pt]
  \small\color{truyengold}--- Ambrose Redmoon ---
\end{center}
\ngancach

\section{Rosa Parks -- Người Phụ Nữ Không Chịu Đứng Dậy}
\begin{truyen}{Rosa Parks -- Người Phụ Nữ Không Chịu Đứng Dậy}{Hoa Kỳ, 1955}
\chuhoa{N}{}N gày 1 tháng 12 năm 1955, Rosa Parks ngồi trên xe buýt ở Montgomery, Alabama, mệt mỏi sau ngày làm việc dài.\\
\textit{(On December 1, 1955, Rosa Parks sat on a bus in Montgomery, Alabama, tired after a long day of work.)}

Tài xế yêu cầu bà nhường ghế cho một hành khách da trắng -- đúng theo luật phân biệt chủng tộc thời đó.\\
\textit{(The driver asked her to give up her seat to a white passenger -- as required by the segregation laws of the time.)}

Rosa Parks nhìn thẳng về phía trước và nói: "Không."\\
\textit{(Rosa Parks looked straight ahead and said: "No.")}

Bà bị bắt. Nhưng hành động dũng cảm nhỏ bé đó đã châm ngòi cho phong trào tẩy chay xe buýt Montgomery.\\
\textit{(She was arrested. But that small act of courage ignited the Montgomery Bus Boycott.)}

Phong trào đó trở thành bước ngoặt của phong trào quyền dân sự, thay đổi lịch sử nước Mỹ mãi mãi.\\
\textit{(That movement became a turning point of the Civil Rights Movement, changing American history forever.)}

\end{truyen}
\begin{baihoc}
  Đôi khi dũng cảm là một từ "Không" đơn giản -- nói vào đúng lúc, bởi đúng người, vì điều đúng đắn.\\
  \textit{(Sometimes courage is one simple word "No" -- said at the right moment, by the right person, for the right reason.)}
\end{baihoc}

\section{Malala Yousafzai -- Đứng Dậy Sau Khi Bị Bắn}
\begin{truyen}{Malala Yousafzai -- Đứng Dậy Sau Khi Bị Bắn}{Pakistan, Thế Kỷ 21}
\chuhoa{M}{}M alala Yousafzai lớn lên ở thung lũng Swat, Pakistan, nơi Taliban cấm phụ nữ đến trường học.\\
\textit{(Malala Yousafzai grew up in Pakistan's Swat Valley, where the Taliban banned women from attending school.)}

Năm 15 tuổi, cô viết blog ẩn danh cho BBC về nỗi sợ hãi và mong muốn được học của cô.\\
\textit{(At age 15, she wrote an anonymous blog for the BBC about her fears and desire to learn.)}

Tháng 10 năm 2012, một tay súng Taliban leo lên xe buýt học sinh và bắn Malala vào đầu.\\
\textit{(In October 2012, a Taliban gunman boarded a school bus and shot Malala in the head.)}

Sau nhiều tháng hồi phục, cô trở lại và tiếp tục mạnh mẽ hơn: "Họ muốn im lặng tôi -- họ đã thất bại."\\
\textit{(After months of recovery, she returned and continued stronger: "They wanted to silence me -- they failed.")}

Năm 2014, ở tuổi 17, Malala trở thành người trẻ nhất nhận giải Nobel Hòa Bình.\\
\textit{(In 2014, at age 17, Malala became the youngest person to receive the Nobel Peace Prize.)}

\end{truyen}
\begin{baihoc}
  Dũng cảm không có nghĩa là không bị thương; nó có nghĩa là tiếp tục đứng dậy dù đã bị thương.\\
  \textit{(Courage does not mean not being wounded; it means continuing to stand up even after being wounded.)}
\end{baihoc}

\section{Gandhi -- Muối Và Tự Do}
\begin{truyen}{Gandhi -- Muối Và Tự Do}{Ấn Độ, 1930}
\chuhoa{N}{}N gày 12 tháng 3 năm 1930, Gandhi bắt đầu cuộc hành trình 390 km đến bờ biển Dandi.\\
\textit{(On March 12, 1930, Gandhi began a 390-kilometer journey to the Dandi seashore.)}

Ông đi để phản đối luật của người Anh độc quyền bán muối -- một thứ thiết yếu nhất cho người nghèo Ấn Độ.\\
\textit{(He walked to protest the British salt monopoly law -- the most essential commodity for India's poor.)}

Hai mươi bốn ngày sau, ông quỳ xuống bờ biển và nhặt một nhúm muối từ đất.\\
\textit{(Twenty-four days later, he knelt at the shoreline and picked up a pinch of salt from the ground.)}

Hành động đơn giản đó phá vỡ luật của đế quốc Anh -- và cả thế giới dõi mắt theo.\\
\textit{(That simple act broke the law of the British Empire -- and the whole world watched.)}

Cuộc hành trình Muối của Gandhi chứng minh: dũng cảm không phải bom đạn mà là sự thật được hành động hóa.\\
\textit{(Gandhi's Salt March proved: courage is not bombs but truth made into action.)}

\end{truyen}
\begin{baihoc}
  Dũng cảm thật sự không cần bạo lực; đôi khi một hành động hòa bình nhỏ có thể làm rung chuyển đế quốc.\\
  \textit{(True courage needs no violence; sometimes a small peaceful act can shake an empire.)}
\end{baihoc}

\section{Nguyễn Trãi -- Dũng Cảm Viết Sự Thật}
\begin{truyen}{Nguyễn Trãi -- Dũng Cảm Viết Sự Thật}{Việt Nam, Thế Kỷ 15}
\chuhoa{K}{}K hi viết Bình Ngô Đại Cáo sau chiến thắng đuổi giặc Minh, Nguyễn Trãi không chỉ tuyên ngôn chiến thắng.\\
\textit{(When writing the Proclamation of Victory after expelling the Ming invaders, Nguyen Trai did not merely declare triumph.)}

Ông tuyên bố nguyên tắc nhân nghĩa: "Việc nhân nghĩa cốt ở yên dân; quân điếu phạt trước lo trừ bạo."\\
\textit{(He proclaimed the principle of benevolence: "The cause of benevolence lies in protecting the people; righteous warfare aims first at removing tyranny.")}

Ông dũng cảm tuyên bố độc lập văn hóa Việt Nam khi cả vùng vẫn trong vòng ảnh hưởng phương Bắc.\\
\textit{(He boldly proclaimed Vietnamese cultural independence when the region was still under Northern influence.)}

"Núi sông bờ cõi đã chia, phong tục Bắc Nam cũng khác." -- Một tuyên ngôn dũng cảm vượt thời đại.\\
\textit{("Mountains and rivers are divided, the customs of North and South also differ." -- A courageous declaration ahead of its time.)}

Dù sau này bị oan khuất, Bình Ngô Đại Cáo vẫn là bản tuyên ngôn độc lập đẹp nhất trong lịch sử Việt Nam.\\
\textit{(Though later wrongly persecuted, the Proclamation remains the most beautiful declaration of independence in Vietnamese history.)}

\end{truyen}
\begin{baihoc}
  Dũng cảm của người trí thức là dám viết sự thật lịch sử dù biết điều đó có thể gây nguy hiểm cho bản thân.\\
  \textit{(The courage of an intellectual is daring to write historical truth knowing it may endanger oneself.)}
\end{baihoc}

\section{Mandela -- 27 Năm Không Khuất Phục}
\begin{truyen}{Mandela -- 27 Năm Không Khuất Phục}{Nam Phi, Thế Kỷ 20}
\chuhoa{T}{}T rong 27 năm tù, chính phủ Nam Phi nhiều lần đề nghị thả Mandela với một điều kiện duy nhất:\\
\textit{(During 27 years in prison, the South African government repeatedly offered to free Mandela with one condition:)}

Ông phải công khai từ bỏ đấu tranh vũ trang và hứa không chống đối chính quyền nữa.\\
\textit{(He must publicly renounce armed struggle and promise not to oppose the government anymore.)}

Mandela từ chối mỗi lần, dù điều đó có nghĩa là tiếp tục giam cầm, lao động khổ sai, cô lập.\\
\textit{(Mandela refused each time, even though it meant continued imprisonment, hard labor, isolation.)}

Ông nói với con gái qua song sắt: "Chỉ người tự do mới có thể đàm phán; tù nhân không thể ký kết.''\\
\textit{(He told his daughter through prison bars: "Only free men can negotiate; prisoners cannot enter contracts.")}

Năm 1990 ông bước ra tự do không điều kiện -- vì cả thế giới đã đứng về phía ông.\\
\textit{(In 1990 he walked out to unconditional freedom -- because the whole world had stood on his side.)}

\end{truyen}
\begin{baihoc}
  Dũng cảm thật sự là giữ nguyên tắc của mình ngay cả khi từ bỏ chúng sẽ mang lại tự do tức thì.\\
  \textit{(True courage is maintaining your principles even when abandoning them would bring immediate freedom.)}
\end{baihoc}

\section{Người Thầy Dám Dạy Sự Thật}
\begin{truyen}{Người Thầy Dám Dạy Sự Thật}{Việt Nam, Hiện Đại}
\chuhoa{C}{}C ó người thầy dạy lịch sử ở trường làng nhỏ, không nổi tiếng, không sách giáo khoa bóng bẩy.\\
\textit{(There was a history teacher at a small village school, not famous, with no glossy textbooks.)}

Mỗi giờ học, ông không chỉ đọc nguyên văn sách giáo khoa -- ông kể những câu chuyện thật, cả phần đau.\\
\textit{(Each lesson, he did not just read the textbook verbatim -- he told real stories, including the painful parts.)}

Ông dạy học trò không chỉ nhớ ngày tháng và chiến thắng, mà nhớ con người và bài học nhân sinh.\\
\textit{(He taught students not just to remember dates and victories, but to remember people and life lessons.)}

Nhiều năm sau, một học trò cũ trở thành giáo viên cũng dạy theo cách đó. Rồi học trò của học trò cũng vậy.\\
\textit{(Years later, a former student who became a teacher also taught that way. Then their students' students too.)}

Dũng cảm của người thầy không cần sân khấu lớn -- nó thể hiện trong từng giờ dạy thật mỗi ngày.\\
\textit{(A teacher's courage needs no grand stage -- it shows in each honest lesson taught every day.)}

\end{truyen}
\begin{baihoc}
  Dũng cảm bình thường của người bình thường -- làm đúng điều mình tin trong công việc hàng ngày -- thay đổi thế giới.\\
  \textit{(The ordinary courage of ordinary people -- doing right by their beliefs in daily work -- changes the world.)}
\end{baihoc}

\section{Trần Hưng Đạo Và Lời Thề Sông Hóa}
\begin{truyen}{Trần Hưng Đạo Và Lời Thề Sông Hóa}{Việt Nam, Thế Kỷ 13}
\chuhoa{N}{}N ăm 1285, quân Nguyên -- Mông ồ ạt xâm lược với 500,000 quân, đội quân đông nhất thế giới lúc đó.\\
\textit{(In 1285, the Mongol Yuan army invaded massively with 500,000 troops, the world's largest army at the time.)}

Trần Hưng Đạo chỉ huy quân dân Đại Việt với số lượng ít hơn nhiều lần nhưng không hề nao núng.\\
\textit{(Tran Hung Dao commanded the Dai Viet forces, far outnumbered, yet without the slightest wavering.)}

Tướng quân thề trước quân sĩ và dòng sông: "Bước qua xác ta mới sang được bờ này."\\
\textit{(The general vowed before his troops and the river: "Only over my body shall you cross to this shore.")}

Câu nói đó không phải hùng biện -- đó là lời cam kết của một con người đã chấp nhận cái chết để bảo vệ đất nước.\\
\textit{(Those words were not rhetoric -- they were the pledge of a man who had accepted death to protect his country.)}

Ba lần Nguyên Mông xâm lược, ba lần bị đánh bại -- lịch sử ghi nhận một trong những chiến thắng kỳ diệu nhất.\\
\textit{(Three times the Mongols invaded, three times they were repelled -- history records one of the most miraculous victories.)}

\end{truyen}
\begin{baihoc}
  Dũng cảm của người lãnh đạo sáng suốt không chỉ ở sức mạnh cá nhân mà ở khả năng truyền lửa tin tưởng cho người khác.\\
  \textit{(The courage of a wise leader lies not only in personal strength but in the ability to ignite belief in others.)}
\end{baihoc}
